\chapter{Orientação, Apresentação e Avaliação}

A avaliação do TCC considera mérito acadêmico, coerência metodológica, qualidade da escrita, normalização e domínio do estudante sobre o próprio trabalho. Cada curso pode definir critérios e formas de apresentação específicos.

\section{Papel do orientador}

O orientador acompanha o desenvolvimento do TCC, mas não é coautor automático do texto. Sua função é ajudar o estudante a qualificar decisões acadêmicas, corrigir rumos, indicar bibliografia, avaliar viabilidade e preparar o trabalho para avaliação.

O estudante deve respeitar os prazos combinados e enviar versões com antecedência. Orientações produtivas dependem de material concreto para leitura.

\section{Papel da banca}

A banca examina o trabalho, identifica contribuições, aponta fragilidades e avalia se o estudante atingiu os objetivos esperados. Entre os aspectos avaliáveis estão:

\begin{itemize}
  \item relevância e delimitação do tema;
  \item clareza do problema e dos objetivos;
  \item adequação metodológica e ética;
  \item qualidade da fundamentação e da análise;
  \item coerência da parte conclusiva;
  \item uso adequado de citações e referências;
  \item apresentação formal e domínio do estudante.
\end{itemize}

\section{Apresentação do trabalho}

A apresentação comunica o TCC já produzido; ela não constitui uma modalidade do documento. O curso ou o evento pode exigir defesa oral, slides, pôster ou combinação dessas formas. Tempo, dimensões, arquivo, identidade visual e critérios de avaliação devem seguir a convocação específica.

\subsection{Defesa oral e slides}

Na defesa oral, o estudante deve expor tema, problema, objetivos, método, resultados, conclusão e contribuições dentro do tempo definido. Os slides devem apoiar a fala com títulos informativos, texto enxuto e elementos visuais legíveis. A apresentação não deve tentar reproduzir todas as páginas do TCC.

É recomendável ensaiar com o arquivo final, testar fontes, vídeos e conexões, manter uma cópia em formato PDF e conferir previamente os recursos disponíveis na sala.

\subsection{Pôster técnico ou científico}

O pôster é uma forma de comunicação visual sintética, comum em seminários e congressos. Deve permitir leitura rápida e apoiar a explicação presencial do estudante. Em geral, apresenta:

\begin{itemize}
  \item título, autoria, orientação e vínculo institucional;
  \item contexto ou introdução breve;
  \item objetivo;
  \item método;
  \item resultados principais;
  \item conclusão;
  \item referências essenciais e informações de contato, quando solicitadas.
\end{itemize}

Gráficos, tabelas, fluxos e imagens devem ter função informativa, fonte identificada e resolução suficiente. Excesso de texto, fonte pequena e baixo contraste prejudicam a leitura. O modelo e as dimensões do evento prevalecem sobre recomendações genéricas.

\section{Arguição}

Durante a arguição, o estudante deve ouvir atentamente, anotar recomendações e responder com clareza. Quando não souber responder, é melhor reconhecer a limitação do que improvisar justificativas frágeis.

Discordâncias acadêmicas podem ocorrer. O estudante pode defender suas escolhas com respeito, base argumentativa e abertura à revisão.

\section{Correções após a banca}

Quando a banca indicar correções, o estudante deve incorporá-las na versão final dentro do prazo estabelecido. A versão completa deve ser revisada novamente, pois alterações pontuais podem afetar sumário, paginação, referências e coerência textual.
